\begingroup
\renewcommand{\arraystretch}{1.45}
\newcommand{\TransscendentArticle}[1]{\par\medskip\Needspace{5\baselineskip}\pdfbookmark[1]{Art. #1}{transscendentes-ii-art-#1}\centerline{[#1.]}\nobreak\smallskip}
\fancyhead[L]{\small ZUR THEORIE DER NEUEN TRANSZENDENTEN. II.}

\Needspace{20\baselineskip}
\TransscendentArticle{8}
Zum Beweis der schönen Lehrsätze der Reziprozität wird Folgendes dienen:
I.\quad Das Produkt aus allen
\[
1-\frac{\eta}{ak+N}
\]
ist das Produkt aus allen
\SourceLeaf{452}
\[
\frac{1-\dfrac{\eta-N}{ak}}{1+\dfrac{N}{ak}}
\]
\noindent also
\[
\begin{gathered}=\frac{\sin\dfrac{N-\eta}{a}\pi}{\sin\dfrac{N}{a}\pi}\\[9pt]=e^{-\frac{\eta\pi i}{a}}\mathbin{.}
 \frac{1-e^{-\frac{2(N-\eta)}{a}\pi i}}{1-e^{-\frac{2N}{a}\pi i}}\\[9pt]=e^{\frac{\eta\pi i}{a}}\mathbin{.}
 \frac{1-e^{\frac{2(N-\eta)}{a}\pi i}}{1-e^{\frac{2N}{a}\pi i}}\end{gathered}
\]

\Needspace{19\baselineskip}
II. Sollen aber nur für $k$ die ungeraden Zahlen eingesetzt werden, so ist
jenes Produkt
\[
\begin{gathered}=\frac{\cos\dfrac{N-\eta}{a}\dfrac{\pi}{2}}{\cos\dfrac{N}{a}\dfrac{\pi}{2}}\\[9pt]=e^{-\frac{\eta\pi i}{2a}}\mathbin{.}
 \frac{1+e^{-\frac{N-\eta}{a}\pi i}}{1+e^{-\frac{N}{a}\pi i}}\\[9pt]=e^{\frac{\eta\pi i}{2a}}\mathbin{.}
 \frac{1+e^{\frac{N-\eta}{a}\pi i}}{1+e^{\frac{N}{a}\pi i}}\end{gathered}
\]
\Needspace{18\baselineskip}
Setzt man also $N=k'\alpha'$ und
\[
e^{-\frac{\alpha'\pi i}{a}}=x\quad\text{und}\quad
e^{\frac{\eta\pi i}{a}}=y
\]
so ergibt sich für den Wert jenes Produkts
\[
=y^{-\frac12}\mathbin{.}\frac{1+x^k y}{1+x^k}
=y^{\frac12}\mathbin{.}\frac{1+x^{-k}\mathbin{.}\dfrac1y}{1+x^{-k}}
\]
\Needspace{8\baselineskip}
Folglich ist das Produkt aus allen Werten, wenn für $k'$ alle ungeraden Zahlen eingesetzt werden,
\[
=(x,y)\frac{[x]^2[x^4]^2}{[xx]^4}
\]
\Needspace{15\baselineskip}
Es seien $m,n$ zwei beliebige positive reelle Größen und $\lambda$ reell oder
imaginär. Man setze
\[
e^{-\frac mn\pi}=x,\qquad e^{\frac\lambda n\pi}=y
\]
\Needspace{13\baselineskip}
Dann ist das Produkt aus allen
\[
1-\frac{\lambda}{km+k'ni}
\]
für $k$ und $k'$ alle ungeraden ganzen Zahlen eingesetzt, dem Obigen zufolge

\SourceLeaf{453}
\[
=(x,y)\frac{[x]^2[x^4]^2}{[xx]^4}
\]
\Needspace{9\baselineskip}
\noindent Offenbar ist aber obiges Produkt auch das Produkt aus allen
\[
1-\frac{\lambda i}{kmi+k'n}
\]
\Needspace{12\baselineskip}
Setzt man also
\[
e^{-\frac nm\pi}=x',\qquad e^{\frac\lambda m\pi i}=y'
\]
so ist jenes Produkt
\[
=(x',y')\frac{[x']^2[x'^4]^2}{[x'x']^4}
\]
\Needspace{19\baselineskip}
Folglich sind diese beiden Ausdrücke einander gleich, oder
\[
\begin{gathered}\frac{1+xy\mathbin{.}1+\dfrac xy\mathbin{.}1+x^3y\mathbin{.}1+\dfrac{x^3}y\mathinner{.\,.\,.}}
 {1+x\mathbin{.}1+x\mathbin{.}1+x^3\mathbin{.}1+x^3\mathinner{.\,.\,.}}\\[9pt]=
\frac{1+x'y'\mathbin{.}1+\dfrac{x'}{y'}\mathbin{.}1+x'^3y'\mathbin{.}1+\dfrac{x'^3}{y'}\mathinner{.\,.\,.}}
 {1+x'\mathbin{.}1+x'\mathbin{.}1+x'^3\mathbin{.}1+x'^3\mathinner{.\,.\,.}}\end{gathered}
\]
Diese Schlüsse bedürfen einer Verbesserung (alle geradezu noch einen konstanten Teil im Nenner).


\Needspace{19\baselineskip}
\TransscendentArticle{9}
\[
\begin{gathered}
\frac{x+4x^4+9x^9+16x^{16}+\mathinner{.\,.}}{1+2x+2x^4+2x^9+\mathinner{.\,.}}
\\[9pt]=\frac{x}{1-xx}-\frac{2xx}{1-x^4}+\frac{3x^3}{1-x^6}-\mathinner{.\,.\,.}
\qquad \mathrm{I}\\[9pt]=\frac{x}{(1+x)^2}+\frac{x^3}{(1+x^3)^2}+\frac{x^5}{(1+x^5)^2}+\mathinner{.\,.\,.}
\qquad \mathrm{II}
\end{gathered}
\]
also
\[
\begin{gathered}\tfrac12\log(1+2x+2x^4+2x^9+\mathinner{.\,.})\\[9pt]=\frac{x}{1+x}+\frac{x^3}{3(1+x^3)}+\frac{x^5}{5(1+x^5)}+\mathinner{.\,.\,.}\end{gathered}
\]
\Needspace{20\baselineskip}
I. Aus der Differentiation des Logarithmus des Ausdrucks in Produktform.

Wird I. entwickelt in
\[
\begin{array}{rrrrl}
x&+x^3&+x^5&+x^7&+\mathinner{.\,.\,.}\\
-2xx&-2x^6&-2x^{10}&-2x^{14}&-\mathinner{.\,.\,.}\\
+3x^3&+3x^9&+3x^{15}&+3x^{21}&+\mathinner{.\,.\,.}\\
-4x^4&-4x^{12}&-4x^{20}&-4x^{28}&-\mathinner{.\,.\,.}\\
\multicolumn{5}{c}{.\quad .\quad .\quad .\quad .\quad .}
\end{array}
\]
und werden die senkrechten Reihen einzeln summiert, so entsteht II.


\Needspace{13\baselineskip}
\SourceLeaf{454}
\[
\begin{gathered}(1+2x+2x^4+2x^9+\mathinner{.\,.})^4\\[7pt]=1+\frac{8x}{1-x}+\frac{16xx}{1+xx}+\frac{24x^3}{1-x^3}+\frac{32x^4}{1+x^4}+\mathinner{.\,.\,.}\end{gathered}
\]
\[
\begin{gathered}(1-2x+2x^4-2x^9+\mathinner{.\,.})^4\\[7pt]=1-\frac{8x}{1+x}+\frac{16xx}{1+xx}-\frac{24x^3}{1+x^3}+\frac{32x^4}{1+x^4}-\mathinner{.\,.\,.}\end{gathered}
\]
\[
\begin{gathered}(2x^{\frac14}+2x^{\frac94}+2x^{\frac{25}4}+\mathinner{.\,.})^4\\[7pt]=\frac{16x}{1-xx}+\frac{48x^3}{1-x^6}+\frac{80x^5}{1-x^{10}}+\mathinner{.\,.\,.}\end{gathered}
\]
\[
\begin{gathered}(1+2x+2x^4+2x^9+\mathinner{.\,.})^4\\[7pt]=1+\frac{8x}{(1-x)^2}+\frac{8xx}{(1+xx)^2}+\frac{8x^3}{(1-x^3)^2}+\frac{8x^4}{(1+x^4)^2}+\mathinner{.\,.\,.}\end{gathered}
\]
\[
\begin{gathered}(1-2x+2x^4-2x^9+\mathinner{.\,.})^4\\[7pt]=1-\frac{8x}{(1+x)^2}+\frac{8xx}{(1+xx)^2}-\frac{8x^3}{(1+x^3)^2}+\frac{8x^4}{(1+x^4)^2}-\mathinner{.\,.}\end{gathered}
\]
\[
\begin{gathered}(2x^{\frac14}+2x^{\frac94}+2x^{\frac{25}4}+\mathinner{.\,.})^4\\[7pt]=\frac{16(1+xx)x}{(1-xx)^2}+\frac{16(1+x^6)x^3}{(1-x^6)^2}+\frac{16(1+x^{10})x^5}{(1-x^{10})^2}+\mathinner{.\,.\,.}\end{gathered}
\]
\Needspace{27\baselineskip}
Die Reihen
\[
\begin{aligned}
p&=1+2x+2x^4+\mathrm{usw.},&\quad\frac1{pp}&=t\\
q&=1-2x+2x^4-\mathrm{usw.},&\quad\frac1{qq}&=u
\end{aligned}
\]
werden durch Differentialgleichungen am einfachsten auf folgende Art ausgedrückt
\[
\begin{array}{lll}
\dfrac{x\,\mathrm{d}t}{\mathrm{d}x}=t',&
\dfrac{x\,\mathrm{d}t'}{\mathrm{d}x}=t'',&
\dfrac{x\,\mathrm{d}t''}{\mathrm{d}x}=t'''\\[9pt]
\dfrac{x\,\mathrm{d}u}{\mathrm{d}x}=u',&
\dfrac{x\,\mathrm{d}u'}{\mathrm{d}x}=u'',&
\dfrac{x\,\mathrm{d}u''}{\mathrm{d}x}=u'''
\end{array}
\]
\[
\frac ut-\frac tu=2(tu'-ut')=-4u^3t''=+4t^3u''
\]
\[
\frac{t'''}{t''}+3\frac{t'}t
=\surd\bigl(\frac1{t^4}+16\frac{t''}t\bigr)
\]
\par\nobreak\medskip\centerline{\rule{.28\textwidth}{.3pt}}
\clearpage
\endgroup
